% GENERATED FILE. Edit canonical lessons, then run npm run generate:books.
% canonical-source-hash: fnv1a64:45c0d3d97f2a2ae2
% canonical-lessons: FR-C465-tempete, FR-C465-inondation, FR-C465-inonder, FR-C465-deborder, FR-C465-flaque, FR-R465-revision-464-465, FR-R465-revision-final-465

\chapter{Tempête et inondations}
\label{ch:fr-tempete-et-inondations}
\hlchaptermodality{465}

\begin{chapteropening}
\textbf{By the end of this chapter:} \emph{I can talk about storms and floods: the storm, the flood, flooding, overflowing and the puddle.}

Complete the last lesson of chapter 465: I can talk about storms and floods: the storm, the flood, flooding, overflowing and the puddle.
\end{chapteropening}

\section[la tempête]{la tempête — the storm, the gale}

\label{lesson:FR-C465-tempete}



\begin{quote}
Before the new one: say the French for clammy, then the French for to sweat.
\end{quote}

\subsection*{You'll want to know: la tempête}
\textbf{la tempête} — \textquotedblleft{}the storm, the gale\textquotedblright{}.

\begin{quote}
La tempête a fait tomber des arbres sur la route.
\end{quote}

\begin{quote}
\emph{The storm brought trees down onto the road.}
\end{quote}

Stronger than an \emph{orage}: a lot of wind.

\begin{cousinweb}
\emph{tempête} comes from Old French \emph{tempeste}, from Vulgar Latin \emph{tempesta}, from Latin \emph{tempestās} 'weather, storm', built on \emph{tempus} 'time' (EtymDB).
\end{cousinweb}

\begin{grammarlens}[title={What you've built}]
One more word for telling, explaining and describing at B1.
\end{grammarlens}

\subsection*{Your turn}
\begin{itemize}
\raggedright
  \item \emph{Say it:} \emph{la tempête}
  \item \emph{Say it:} \emph{la tempête}, once more
  \item \emph{Say it:} a sentence of your own with \emph{la tempête}
  \item \emph{Recall:} say the French for clammy, then the French for to sweat, then say \emph{la tempête} again
\end{itemize}

\begin{tcolorbox}[breakable,colback=teal!4,colframe=teal!35!black,title={Before you move on}]
What does la tempête mean? (\textquotedblleft{}The storm, the gale\textquotedblright{}.) Say it once more.
\end{tcolorbox}
\section[l'inondation]{l'inondation — the flood}

\label{lesson:FR-C465-inondation}



\begin{quote}
Before the new one: say the French for to sweat, then the French for the storm.
\end{quote}

\subsection*{You'll want to know: l'inondation}
\textbf{l'inondation} — \textquotedblleft{}the flood\textquotedblright{}.

\begin{quote}
Après l'inondation, la cave était pleine d'eau.
\end{quote}

\begin{quote}
\emph{After the flood the cellar was full of water.}
\end{quote}

A feminine noun.

\begin{cousinweb}
\emph{inondation} is a learned borrowing from Latin \emph{inundātiō} (EtymDB), from \emph{inundāre} 'to overflow', built on \emph{unda} 'wave'.
\end{cousinweb}

\begin{grammarlens}[title={What you've built}]
One more word for telling, explaining and describing at B1.
\end{grammarlens}

\subsection*{Your turn}
\begin{itemize}
\raggedright
  \item \emph{Say it:} \emph{l'inondation}
  \item \emph{Say it:} \emph{l'inondation}, once more
  \item \emph{Say it:} a sentence of your own with \emph{l'inondation}
  \item \emph{Recall:} say the French for to sweat, then the French for the storm, then say \emph{l'inondation} again
\end{itemize}

\begin{tcolorbox}[breakable,colback=teal!4,colframe=teal!35!black,title={Before you move on}]
What does l'inondation mean? (\textquotedblleft{}The flood\textquotedblright{}.) Say it once more.
\end{tcolorbox}
\section[inonder]{inonder — to flood}

\label{lesson:FR-C465-inonder}



\begin{quote}
Before the new one: say the French for the storm, then the French for the flood.
\end{quote}

\subsection*{You'll want to know: inonder}
\textbf{inonder} — \textquotedblleft{}to flood\textquotedblright{}.

\begin{quote}
La rivière a inondé les champs.
\end{quote}

\begin{quote}
\emph{The river flooded the fields.}
\end{quote}

Figurative: \emph{inonder de messages} — to flood with messages.

\begin{cousinweb}
\emph{inonder} is a learned borrowing from Latin \emph{inundāre} 'to overflow' (EtymDB), from \emph{unda} 'wave'.
\end{cousinweb}

\begin{grammarlens}[title={What you've built}]
One more word for telling, explaining and describing at B1.
\end{grammarlens}

\subsection*{Your turn}
\begin{itemize}
\raggedright
  \item \emph{Say it:} \emph{inonder}
  \item \emph{Say it:} \emph{inonder}, once more
  \item \emph{Say it:} a sentence of your own with \emph{inonder}
  \item \emph{Recall:} say the French for the storm, then the French for the flood, then say \emph{inonder} again
\end{itemize}

\begin{tcolorbox}[breakable,colback=teal!4,colframe=teal!35!black,title={Before you move on}]
What does inonder mean? (\textquotedblleft{}To flood\textquotedblright{}.) Say it once more.
\end{tcolorbox}
\section[déborder]{déborder — to overflow; to go beyond}

\label{lesson:FR-C465-deborder}



\begin{quote}
Before the new one: say the French for the flood, then the French for to flood.
\end{quote}

\subsection*{You'll want to know: déborder}
\textbf{déborder} — \textquotedblleft{}to overflow; to go beyond\textquotedblright{}.

\begin{quote}
Après trois jours de pluie, le ruisseau déborde.
\end{quote}

\begin{quote}
\emph{After three days of rain the stream is overflowing.}
\end{quote}

\emph{Être débordé} — to be swamped with work.

\begin{cousinweb}
\emph{déborder} is built from \emph{dé-} and \emph{bord} 'edge, side' (EtymDB lists it), which comes from a Frankish word for the side of a ship.
\end{cousinweb}

\begin{grammarlens}[title={What you've built}]
One more word for telling, explaining and describing at B1.
\end{grammarlens}

\subsection*{Your turn}
\begin{itemize}
\raggedright
  \item \emph{Say it:} \emph{déborder}
  \item \emph{Say it:} \emph{déborder}, once more
  \item \emph{Say it:} a sentence of your own with \emph{déborder}
  \item \emph{Recall:} say the French for the flood, then the French for to flood, then say \emph{déborder} again
\end{itemize}

\begin{tcolorbox}[breakable,colback=teal!4,colframe=teal!35!black,title={Before you move on}]
What does déborder mean? (\textquotedblleft{}To overflow; to go beyond\textquotedblright{}.) Say it once more.
\end{tcolorbox}
\section[la flaque]{la flaque — the puddle}

\label{lesson:FR-C465-flaque}



\begin{quote}
Before the new one: say the French for to flood, then the French for to overflow.
\end{quote}

\subsection*{You'll want to know: la flaque}
\textbf{la flaque} — \textquotedblleft{}the puddle\textquotedblright{}.

\begin{quote}
Le bus est passé dans une flaque et m'a mouillé.
\end{quote}

\begin{quote}
\emph{The bus went through a puddle and got me wet.}
\end{quote}

\emph{Une flaque d'eau}.

\begin{cousinweb}
\emph{flaque} came into French from Picard \emph{flaque} (EtymDB), a northern dialect word usually linked to \emph{flache} 'hollow, soft ground'; the deeper origin is not settled.
\end{cousinweb}

\begin{grammarlens}[title={What you've built}]
One more word for telling, explaining and describing at B1.
\end{grammarlens}

\subsection*{Your turn}
\begin{itemize}
\raggedright
  \item \emph{Say it:} \emph{la flaque}
  \item \emph{Say it:} \emph{la flaque}, once more
  \item \emph{Say it:} a sentence of your own with \emph{la flaque}
  \item \emph{Recall:} say the French for to flood, then the French for to overflow, then say \emph{la flaque} again
\end{itemize}

\begin{tcolorbox}[breakable,colback=teal!4,colframe=teal!35!black,title={Before you move on}]
What does la flaque mean? (\textquotedblleft{}The puddle\textquotedblright{}.) Say it once more.
\end{tcolorbox}
\section[(dialogue)]{Ville, commerces, nature et météo}

\label{lesson:FR-R465-revision-464-465}



\begin{quote}
Say the French for the heatwave and the drought.
\end{quote}

\subsection*{Your turn}
Say these aloud:

\begin{itemize}
\raggedright
  \item the heatwave, the drought, the humidity, clammy, to sweat
  \item the storm, the flood, to flood, to overflow, the puddle
\end{itemize}

\begin{tcolorbox}[breakable,colback=teal!4,colframe=teal!35!black,title={Before you move on}]
The heatwave? (\textbf{la canicule}.) The drought? (\textbf{la sécheresse}.) The humidity? (\textbf{l'humidité}.) Clammy? (\textbf{moite}.) To sweat? (\textbf{transpirer}.) The storm? (\textbf{la tempête}.) The flood? (\textbf{l'inondation}.) To flood? (\textbf{inonder}.) To overflow? (\textbf{déborder}.) The puddle? (\textbf{la flaque}.)
\end{tcolorbox}
\section[(dialogue)]{Tempête et inondations — once more}

\label{lesson:FR-R465-revision-final-465}



\begin{quote}
Say the French for the storm and the flood.
\end{quote}

\subsection*{Your turn}
\begin{itemize}
\raggedright
  \item \emph{Say it:} the storm, the flood, to flood, to overflow, the puddle
\end{itemize}

\begin{tcolorbox}[breakable,colback=teal!4,colframe=teal!35!black,title={Before you move on}]
The storm? (\textbf{la tempête}.) The flood? (\textbf{l'inondation}.) To flood? (\textbf{inonder}.) To overflow? (\textbf{déborder}.) The puddle? (\textbf{la flaque}.)
\end{tcolorbox}
